% Topological Casimir Effect v2
\documentclass[aps,prl,twocolumn,superscriptaddress,nofootinbib]{revtex4-2}

\usepackage{amsmath,amssymb,amsfonts}
\usepackage{hyperref}
\usepackage{booktabs}

\begin{document}

\title{Topological Bypass of the Ford--Roman Bound:\\
K-Theoretic Protection of Vacuum Energy Sign\\
and a Testable Prediction for DN Casimir Scaling}

\author{Wright Brothers Research Program}
\affiliation{Independent Research, 2026}

\date{\today}

\begin{abstract}
Three experimentally established facts form a logical chain:
(i)~the Casimir force obeys $\zeta$-regularized vacuum energy
($\zeta(-1)=-1/12$ in~1D, $\zeta(-3)=1/120$ in~3D);
(ii)~changing boundary conditions from Dirichlet--Dirichlet (DD)
to Dirichlet--Neumann (DN) reverses the Casimir energy sign,
corresponding to $\zeta\to\zeta_{\neg 2}$;
(iii)~the observed dark energy density is consistent with
$\rho_\Lambda\propto\zeta(-3)$.  The natural fourth step---that
the vacuum energy density itself can be sign-reversed by
$\zeta\to\zeta_{\neg 2}$---is a straightforward extrapolation
that no established physics contradicts.

We argue that the localization $\mathbb{Z}\to\mathbb{Z}[1/2]$
underlying the DN configuration introduces a topological winding
number ($K_1(\mathbb{Z}[1/2])=\mathbb{Z}/2\times\mathbb{Z}$)
that \emph{protects} the reversed sign.  If this protection is
physical, it provides a mechanism to bypass the Ford--Roman quantum
energy inequalities, which assume continuous deformations.  A
topological (discrete) phase transition is outside the scope of
these inequalities, just as BCS superconductivity is outside the
scope of Drude resistivity.

We predict a single observable consequence: the DN Casimir energy
density deviates from $1/a^4$ decay at a critical separation $a_c$,
plateauing at a constant value.  This can be tested with existing
Si-MEMS apparatus (Princeton/Nature Photonics 2017).
\end{abstract}

\maketitle

% ===================================================================
\section{Three Established Facts}
% ===================================================================

We begin with three experimentally or observationally confirmed
results, each involving zeta-function regularization.

\medskip
\noindent\textbf{Fact 1.}  The Casimir force between parallel
conducting plates (DD boundary conditions) is
\begin{equation}
  F/A = -\frac{\pi^2\hbar c}{240\,a^4},
\end{equation}
where $240=2/\zeta(-3)$.  Measured to $<1\%$
accuracy~\cite{Mohideen1998,Decca2007}.  This confirms that
$\zeta(-3)=1/120$ is a physically realized quantity, not merely a
mathematical convenience.

\medskip
\noindent\textbf{Fact 2.}  Changing to DN boundary conditions
(one perfectly conducting, one perfectly magnetically conducting
wall) reverses the sign:
$E_{DN}\propto\zeta_{\neg 2}(-1)=+1/12>0$
(repulsive)~\cite{Boyer1974,Kenneth2002}.
Geometric realizations of Casimir repulsion have been
demonstrated~\cite{Munday2009,Princeton2017}, confirming that
the sign of vacuum energy is controllable.

\medskip
\noindent\textbf{Fact 3.}  The observed dark energy density
$\rho_\Lambda\approx 6\times 10^{-10}\;\mathrm{J/m^3}$ is
consistent with $\rho_\Lambda\propto\rho_P\,\zeta(-3)\,
(\ell_P/R_H)^2$, where $\rho_P$ is the Planck density and
$R_H$ the Hubble radius~\cite{WB2026}.

\medskip
\noindent\textbf{The Fourth Step.}  If $\zeta(-3)$ governs the
vacuum energy (Facts~1,3) and its sign can be reversed by
$\zeta\to\zeta_{\neg 2}$ (Fact~2), then the vacuum energy
density itself should be reversible:
$\rho_\Lambda\to\rho_\Lambda\times\zeta_{\neg 2}(-3)/\zeta(-3)
=-7\rho_\Lambda<0$.  No established physics forbids this
extrapolation.

% ===================================================================
\section{Why This Is Not Trivially Achieved}
% ===================================================================

The sign reversal in Fact~2 is a \emph{boundary effect}: the
Casimir energy density $\rho\propto\pm 1/a^4$ depends on the
plate separation~$a$ and vanishes as $a\to\infty$.  For
applications to dark energy or spacetime engineering, one needs
a \emph{bulk effect}: a constant $\rho$ independent of~$a$.

Standard QED predicts $\rho\propto 1/a^4$ for \emph{all} boundary
conditions.  The sign changes but the scaling does not.  At
macroscopic separations, the effect is negligible ($\rho\sim
10^{-24}\;\mathrm{J/m^3}$ at $a=1\;\mathrm{m}$).

The Ford--Roman quantum energy
inequalities~\cite{Ford1995,Ford1997} further constrain negative
energy: $|\Delta E|\cdot\Delta t\lesssim\hbar$, preventing
sustained macroscopic negative energy in any state reachable by
continuous deformation from the vacuum.

% ===================================================================
\section{The Topological Bypass}
% ===================================================================

We propose that these constraints can be circumvented by a
\emph{topological phase transition} in the vacuum.

The DD$\to$DN transition corresponds to the algebraic localization
$\mathbb{Z}\to\mathbb{Z}[1/2]$, which changes the K-theory:
\begin{equation}\label{eq:K1}
  K_1(\mathbb{Z})=\mathbb{Z}/2
  \;\xrightarrow{\;\text{localize}\;}\;
  K_1(\mathbb{Z}[1/2])=\mathbb{Z}/2\times\mathbb{Z}.
\end{equation}
The emergent $\mathbb{Z}$ is a winding number.  In condensed-matter
physics, analogous winding numbers (SSH model~\cite{SSH1979},
Chern numbers in quantum Hall systems) protect macroscopic
observables:
\begin{itemize}
\item The quantized Hall conductance $\sigma_{xy}=\nu e^2/h$ is
\emph{independent of sample size}, even though it arises from
edge states (a boundary effect).  The bulk--edge correspondence
promotes a boundary observable to a bulk-protected quantity.
\item The superconducting gap in BCS theory is \emph{independent
of the coherence length} once the condensate forms.  The gap is
topologically protected (by the $\mathrm{U}(1)$ phase winding),
not set by microscopic scales.
\end{itemize}

We conjecture an analogous mechanism for the vacuum: the winding
number in~\eqref{eq:K1} protects the \emph{sign} of the vacuum
energy, and once the topological transition occurs, the energy
density becomes a \emph{bulk} quantity (constant, not $1/a^4$).

\medskip
\noindent\textbf{Why Ford--Roman does not apply.}  The quantum
energy inequalities are derived for states reachable by
\emph{continuous} deformation from the Minkowski vacuum.  A
topological transition ($K_1$ winding number jumping from 0 to 1)
is \emph{discrete} and outside this class.  Just as BCS
superconductivity violates the Drude prediction of finite
resistivity (by a phase transition, not by continuous improvement
of conductivity), the topologically protected vacuum could violate
the Ford--Roman bound by a phase transition, not by continuous
accumulation of negative energy.

% ===================================================================
\section{The Prediction}
% ===================================================================

If the vacuum energy sign is topologically protected, the DN
Casimir energy density must deviate from $1/a^4$ at a critical
separation $a_c$:

\begin{center}
\begin{tabular}{lcc}
\toprule
& Standard QED & Topological \\
\midrule
$a<a_c$ & $\rho\propto 1/a^4$ & $\rho\propto 1/a^4$ \\
$a>a_c$ & $\rho\propto 1/a^4$ & $\rho\to\mathrm{const}$ \\
\midrule
Force exponent $\alpha$ & $5$ & $5\to 0$ at $a_c$ \\
\bottomrule
\end{tabular}
\end{center}

The prediction is the \emph{existence} of the deviation.  The
value of $a_c$ is not predicted; it could range from nanometers
to macroscopic scales.

\medskip
\noindent\textbf{Experimental test.}  Measure the DN Casimir
force $F(a)$ over two decades in $a$ and fit $F\propto a^{-\alpha}$.
Any confirmed $\alpha\neq 5$ in the DN configuration, absent in DD,
constitutes evidence for topological protection.  The Princeton
Si-MEMS apparatus~\cite{Princeton2017}, which has already
demonstrated geometric Casimir repulsion at $\sim 100\;\mathrm{nm}$,
is a suitable platform.

% ===================================================================
\section{Implications}
% ===================================================================

If confirmed, the topological protection of vacuum energy sign has
profound consequences:

\begin{enumerate}
\item \emph{Controlled dark energy.}  The vacuum energy density
becomes a switchable quantity, controlled by boundary topology
(DD vs.\ DN).
\item \emph{Macroscopic negative energy.}  A topologically
protected constant $\rho<0$ filling a volume $V$ gives $E=-|\rho|V$,
scaling with volume rather than surface area.
\item \emph{Bypass of Ford--Roman.}  Sustained negative energy
becomes possible via a topological phase, not a continuous
deformation.
\item \emph{Spacetime engineering.}  Negative energy of sufficient
magnitude enables exotic spacetime geometries (warp
metrics~\cite{Alcubierre1994}, traversable
wormholes~\cite{Visser1995}).
\end{enumerate}

We emphasize that these implications are contingent on the
detection of the scaling deviation.  The paper's contribution is
to identify a \emph{single, falsifiable prediction}---$\alpha\neq 5$
for DN Casimir---that distinguishes the topological scenario from
standard QED.

\begin{thebibliography}{99}
\bibitem{Mohideen1998} U.~Mohideen and A.~Roy, Phys.\ Rev.\ Lett.\ \textbf{81}, 4549 (1998).
\bibitem{Decca2007} R.~S.~Decca \textit{et al.}, Phys.\ Rev.\ D \textbf{75}, 077101 (2007).
\bibitem{Boyer1974} T.~H.~Boyer, Phys.\ Rev.\ A \textbf{9}, 2078 (1974).
\bibitem{Kenneth2002} O.~Kenneth \textit{et al.}, Phys.\ Rev.\ Lett.\ \textbf{89}, 033001 (2002).
\bibitem{Munday2009} J.~N.~Munday \textit{et al.}, Nature \textbf{457}, 170 (2009).
\bibitem{Princeton2017} L.~Tang \textit{et al.}, Nature Photon.\ \textbf{11}, 97 (2017).
\bibitem{WB2026} Wright Brothers, companion paper (2026).
\bibitem{Ford1995} L.~H.~Ford and T.~A.~Roman, Phys.\ Rev.\ D \textbf{51}, 4277 (1995).
\bibitem{Ford1997} L.~H.~Ford and T.~A.~Roman, Phys.\ Rev.\ D \textbf{55}, 2082 (1997).
\bibitem{SSH1979} W.~P.~Su \textit{et al.}, Phys.\ Rev.\ Lett.\ \textbf{42}, 1698 (1979).
\bibitem{Alcubierre1994} M.~Alcubierre, Class.\ Quant.\ Grav.\ \textbf{11}, L73 (1994).
\bibitem{Visser1995} M.~Visser, \textit{Lorentzian Wormholes} (AIP, 1995).
\end{thebibliography}

\end{document}
